\documentclass[11pt,a4paper]{article}

\usepackage[T1]{fontenc}
\usepackage[utf8]{inputenc}
\usepackage{amsmath}
\usepackage{booktabs}
\usepackage[margin=2.5cm]{geometry}
\usepackage[hidelinks]{hyperref}

\title{The Dirichlet--Multinomial Model and the Choice of Window Length}
\author{LM-IDNet project note}
\date{}

\begin{document}

\maketitle

\section{The important correction}

Each row of the training data is one ten-minute window.  Its four entries are
the packet counts for TCP, UDP, SSDP, and ARP:

\[
\mathbf{x}_r =
\left(
x_{r,\mathrm{TCP}},
x_{r,\mathrm{UDP}},
x_{r,\mathrm{SSDP}},
x_{r,\mathrm{ARP}}
\right).
\]

The production training code does \emph{not} calculate four alpha values for
each row and then average those results.  Instead, it uses all the training
rows to estimate one set of four alpha parameters:

\[
\boldsymbol{\alpha}=
\left(
\alpha_{\mathrm{TCP}},
\alpha_{\mathrm{UDP}},
\alpha_{\mathrm{SSDP}},
\alpha_{\mathrm{ARP}}
\right).
\]

This one set describes normal traffic over the complete training period.  The
estimation question is:

\begin{quote}
Which four alpha values best describe both the usual division of packets among
the four categories and the amount by which that division changes from one
ten-minute window to another?
\end{quote}

There is a separate reporting routine that can estimate a model for an
individual capture or day.  Those daily estimates are descriptive statistics;
they are not averaged to create the production model.

\section{What the four alpha parameters mean}

Define the total concentration

\[
A = \alpha_{\mathrm{TCP}}+
    \alpha_{\mathrm{UDP}}+
    \alpha_{\mathrm{SSDP}}+
    \alpha_{\mathrm{ARP}}.
\]

The expected proportion of category $k$ is

\[
p_k=\frac{\alpha_k}{A}.
\]

The ratios among the four alpha values therefore describe the usual traffic
composition.  For example, a large value of
$\alpha_{\mathrm{UDP}}/A$ means that UDP normally accounts for a large fraction
of the categorized packets.

The sum $A$ has a different meaning.  It describes how tightly ten-minute
windows are expected to remain near the usual composition:

\begin{itemize}
  \item a large $A$ means that normal windows should have fairly similar
        category proportions;
  \item a small $A$ means that normal windows may have substantially different
        category proportions.
\end{itemize}

The project also reports

\[
\psi=\frac{1}{A}.
\]

Consequently, a larger $\psi$ means that the fitted model permits more
window-to-window variation in category proportions.  Summing the four alpha
values to obtain $A$ does not discard the four original values.  The saved
model contains the complete set of alpha values, their sum, their normalized
proportions, and $\psi$.

\section{What is actually summed during estimation}

For a proposed set of alpha values, the program calculates how probable each
ten-minute row is under those values.  Let

\[
P(\mathbf{x}_r\mid\boldsymbol{\alpha})
\]

denote the Dirichlet--Multinomial probability of row $r$.  The current model
assumes that, once the alpha values are fixed, it does not need the preceding
window to calculate the probability of the current window.  Under that
assumption, the probability of the complete collection is

\[
P(\mathbf{x}_1,\ldots,\mathbf{x}_R\mid\boldsymbol{\alpha})
=
\prod_{r=1}^{R}
P(\mathbf{x}_r\mid\boldsymbol{\alpha}).
\]

Products of many small probabilities are inconvenient for a computer because
they can become too close to zero to represent accurately.  Taking logarithms
turns the product into a sum:

\[
\log P(\mathbf{x}_1,\ldots,\mathbf{x}_R\mid\boldsymbol{\alpha})
=
\sum_{r=1}^{R}
\log P(\mathbf{x}_r\mid\boldsymbol{\alpha}).
\]

Thus, the program sums the log-probability contributed by each row.  It does
not sum row-specific alpha estimates, because there are no row-specific alpha
estimates in the production fit.

The estimator repeatedly adjusts the same four alpha values.  After each
adjustment, it recalculates the total log-probability.  It stops when another
adjustment changes that value by less than the configured tolerance, or when it
reaches the configured iteration limit.

\subsection{An initialization detail that can cause confusion}

Before the repeated adjustments begin, the code needs a starting value.  It
adds the counts in each category over the training rows, converts those four
totals into proportions, and multiplies them by a configured starting
concentration.  This produces only the \emph{initial guess}.  It is not the
final model.  The later adjustments use the individual rows and their
window-to-window differences.

\section{Why adding log-probabilities does not pool the windows}

Consider two categories, A and B, and two windows:

\[
\mathbf{x}_1=(90,10),
\qquad
\mathbf{x}_2=(10,90).
\]

The two windows have very different compositions.  If their raw counts were
added before fitting, the result would be

\[
\mathbf{x}_1+\mathbf{x}_2=(100,100).
\]

That single row would show an equal overall composition but would hide the
large change between the two windows.

The current production code does not perform that pooling operation.  It
calculates one log-probability for $(90,10)$, another for $(10,90)$, and then
adds the two log-probabilities.  Because each row is evaluated separately, the
fitted model can learn both of the following facts:

\begin{itemize}
  \item the long-run average is approximately 50 percent A and 50 percent B;
  \item individual windows vary greatly around that average.
\end{itemize}

Therefore, summing the row log-probabilities does not make the
window-to-window variation disappear.  Adding the raw counts first would make
it disappear, but that is not how the production fit works.

\section{What information the current model retains}

The fitted alpha parameters retain information about:

\begin{itemize}
  \item the usual proportion of TCP, UDP, SSDP, and ARP traffic;
  \item how much those proportions differ between normal ten-minute windows;
  \item which individual category compositions are common or unusual under
        the fitted distribution.
\end{itemize}

This is useful when the detection question is:

\begin{quote}
Does the category composition of this ten-minute window resemble the range of
compositions observed during normal training traffic?
\end{quote}

\section{What information the current model discards}

The fitting function receives the count rows but not their timestamps.  It
would produce the same alpha values if the rows were placed in a different
order.  It therefore cannot learn:

\begin{itemize}
  \item that one composition normally occurs in the morning and another at
        night;
  \item that one type of window normally follows another type;
  \item that an upload normally continues for three consecutive windows;
  \item that weekdays and weekends have different patterns;
  \item that separate days have different normal operating conditions.
\end{itemize}

The model also describes how an observed packet total is divided among the
four categories.  It does not provide a separate probability model for why one
window contains 100 packets and another contains 10,000.  An event that changes
traffic volume while preserving the usual category proportions can therefore
be difficult to detect with the composition model alone.

The loss of time information is not caused by adding the log-probabilities.
It is caused by using one time-independent set of alpha parameters and by not
including the order or time of the windows in the model.

\section{Can one time-independent model become too broad?}

Yes.  Suppose normal nighttime traffic is mostly UDP and normal daytime traffic
is mostly TCP.  One time-independent model must explain both regimes with the
same four alpha values.  It may respond by allowing a very wide range of
compositions.

That can cause two problems:

\begin{enumerate}
  \item A genuinely unusual composition may not be flagged because the fitted
        model has become broad enough to tolerate it.
  \item If one normal regime dominates the training data, valid windows from a
        less common regime may be flagged as anomalous.
\end{enumerate}

Thus, the current approach does preserve variation between rows, but it may
combine several distinct daily operating regimes into one overly general
description of normal traffic.

\section{Alternatives for daily and sequential traffic patterns}

No single extension solves every temporal problem.  The appropriate choice
depends on which omitted pattern is present in the data.

\begin{table}[htbp]
\centering
\small
\begin{tabular}{p{0.25\textwidth}p{0.38\textwidth}p{0.27\textwidth}}
\toprule
Method & What it represents & Main limitation \\
\midrule
Broad time-of-day models
& Estimate separate alpha parameters for periods such as night, morning,
  afternoon, and evening.
& Requires enough normal training days in every period. \\
\addlinespace
Time-of-day thresholds
& Keep one alpha model, but compare scores with a threshold appropriate for
  the current period of the day.
& Does not correct a genuinely different expected category composition. \\
\addlinespace
Several Dirichlet--Multinomial modes
& Represent distinct normal activities such as idle, discovery, streaming,
  and cloud upload.
& Represents several activities but does not say which activity should follow
  another. \\
\addlinespace
State-transition model
& Give each operating state its own alpha parameters and estimate which states
  normally follow each other.
& More parameters, more implementation work, and more training data. \\
\addlinespace
Day-level hierarchical model
& Estimate day-specific behavior while sharing information across days.
& More difficult to fit, explain, and validate. \\
\addlinespace
Separate traffic-volume detector
& Detect unusually high or low packet totals even when category proportions
  remain normal.
& Adds a second score that must be calibrated and combined with the existing
  score. \\
\bottomrule
\end{tabular}
\end{table}

If the problem is only that attacks persist for several adjacent windows, a
smaller change is to monitor consecutive anomaly scores or a running average of
recent scores.  This can detect a sustained moderate change without replacing
the Dirichlet--Multinomial model.

\section{Recommended next step for modelling time-of-day behavior}

The present model should remain the comparison baseline.  Before replacing it,
normal training and calibration data should be examined by time of day:

\begin{enumerate}
  \item Plot each category's proportion against the ten-minute position within
        the day, using several normal days.
  \item Plot anomaly scores against time of day and check whether the same
        periods repeatedly have different score distributions.
  \item Check whether adjacent scores remain similar for long runs and whether
        the pattern repeats one day later.
  \item Compare the current model with a small number of broad time-of-day
        models using chronological validation: train on earlier days and test
        on later days.
\end{enumerate}

If a stable daily schedule is visible and the broad-period models improve
later-day results, separate models for night, morning, afternoon, and evening
are the clearest first extension.  Creating one model for every ten-minute
position would require much more data and should only be considered if the
broader periods are demonstrably insufficient.

A separate traffic-volume detector should be evaluated independently because
category composition and packet volume answer different questions.  A more
complex state-transition or hierarchical model is justified only if the
simpler alternatives fail in chronological tests.

\section{Why the project uses ten-minute windows}

Ten minutes is not a mathematically correct window length that can be derived
from the Dirichlet--Multinomial distribution.  Window length is a design
choice.  The best length depends on how quickly the device changes, how long
the attacks last, how many packets the device produces, and how much normal
data is available for fitting and calibration.

The present choice has a clear historical reason: the reference IoT study used
ten-minute windows for the D-Link camera.  Retaining that duration gave this
project a reproducible baseline and allowed direct comparison with the
reference method.  It should therefore be described as the justified
\emph{baseline}, not as a window length that has already been shown to be
optimal for every device or attack.

The current boundary rule is also precise.  Windows are aligned to UTC epoch
boundaries and are closed on the left and open on the right.  For example,

\[
[10{:}00,10{:}10)
\]

contains a packet at 10:00 but not a packet at 10:10.  The packet at 10:10
belongs to the next window.  This prevents a boundary packet from being counted
twice.

\section{Advantages of ten-minute windows}

\subsection{A reasonable number of packets per observation}

Very short windows can contain few packets or only one active category.  Such
rows give little information about the division of traffic among TCP, UDP,
SSDP, and ARP.  Ten minutes gives many IoT devices enough time to produce a
more informative count row while remaining short enough to issue several
decisions per hour.

This does not mean that every ten-minute row must contain every category.  It
means that the complete fitting period is more likely to contain enough
evidence about all four categories and about their normal variation.

\subsection{Moderate resistance to packet-level randomness}

The arrival time of an individual packet can move slightly because of network
delay, retransmission, or scheduling.  In a very short window, a few such
packets can change the observed proportions substantially.  A ten-minute sum
is less sensitive to those small timing changes.

\subsection{Useful operational resolution}

A non-overlapping ten-minute design produces 144 possible decisions in a full
day.  It is much more informative than one model input per hour or one per day,
but its storage and computation costs remain small.

\subsection{Compatibility with the reference study}

Keeping the same duration as the reference study separates two questions:

\begin{enumerate}
  \item whether this project implemented and extended the reference method
        correctly; and
  \item whether another window duration improves anomaly detection.
\end{enumerate}

Changing the duration before establishing the baseline would make those two
effects difficult to distinguish.

\section{Disadvantages of ten-minute windows}

\subsection{Detection is delayed}

A decision cannot normally be made until the window closes.  An attack that
begins just after a boundary may therefore continue for almost ten minutes
before its first complete window is scored.  The average waiting time caused
only by the boundary is approximately five minutes, although processing and
reporting can add further delay.

\subsection{Short attacks can be diluted}

Suppose an attack lasts for 30 seconds and the remaining nine and a half
minutes contain normal traffic.  The resulting row combines both behaviors.
The normal counts may dominate, causing the attack-related change in category
proportions to appear small.

The same issue affects evaluation labels.  The current label is positive when
an attack overlaps any part of the window.  A ten-minute row may therefore be
called an attack row even when only a small fraction of it contains attack
traffic.  A longer window makes this mismatch more severe.

\subsection{Boundaries can split one event}

An event from 10:09 to 10:11 is divided between two rows.  Neither row contains
the complete event, and two weak changes may be harder to detect than one
concentrated change.  Moving the window boundary would produce a different
division even though the underlying packets had not changed.

\subsection{Long activities create related rows}

A cloud upload, scan, or attack may continue across several ten-minute rows.
Those rows contain different packets, but they arise from the same activity and
can be very similar.  The fitting calculation treats their probability
contributions as separate.  This is acceptable for constructing a practical
score, but the rows should not be described as completely independent pieces
of experimental evidence.

\subsection{One duration is unlikely to suit every device}

A busy camera and a nearly silent sensor can produce very different packet
counts in ten minutes.  Ten minutes may be too coarse for a busy device but too
short and sparse for a quiet one.  Attack duration also matters: a rapid scan
and a several-hour exfiltration event should not be expected to have the same
best observation length.

\subsection{The duration affects the statistical model}

Changing the duration changes more than the number of rows.  It changes:

\begin{itemize}
  \item the counts and proportions in each row;
  \item the amount of variation observed between rows;
  \item the fitted alpha parameters and concentration;
  \item the distribution of anomaly scores;
  \item the calibrated threshold;
  \item the number and meaning of positive label rows;
  \item the elapsed time represented by adaptation settings expressed as a
        number of windows.
\end{itemize}

For example, an adaptation interval of 288 windows represents 48 hours when
windows last ten minutes, but 96 hours when they last twenty minutes.  Such
settings must be converted according to elapsed time if different window
durations are to be compared fairly.

\section{The basic trade-off}

\begin{table}[htbp]
\centering
\small
\begin{tabular}{p{0.28\textwidth}p{0.30\textwidth}p{0.30\textwidth}}
\toprule
Property & Shorter windows & Longer windows \\
\midrule
Detection delay
& Faster first decision.
& Slower first decision. \\
\addlinespace
Short attacks
& Less likely to be mixed with many minutes of normal traffic.
& More likely to be diluted by normal traffic. \\
\addlinespace
Count stability
& Sparser and more affected by a few packets.
& Larger, usually more stable counts. \\
\addlinespace
Number of observations
& More rows for fitting and calibration.
& Fewer rows, which is especially harmful for estimating a low score quantile.
  \\
\addlinespace
Sustained changes
& Spread across several related rows.
& Concentrated into fewer rows, but start and end times are less precise. \\
\addlinespace
Operating regimes
& Less mixing of different activities inside one row.
& Greater chance of combining idle, active, and attack behavior in one row. \\
\bottomrule
\end{tabular}
\end{table}

Neither column is uniformly better.  Window duration must be selected using a
predefined development experiment and then left unchanged for the locked final
test.

\section{What can be changed without reading the PCAP files again}

The saved processed JSON files contain, for every ten-minute interval:

\begin{itemize}
  \item the start and end timestamps;
  \item the four category counts;
  \item whether the interval was observed, observed but silent, or missing.
\end{itemize}

Packet counts are additive.  Therefore, a longer window whose length is an
integer multiple of ten minutes can be constructed exactly from adjacent saved
rows.  For example,

\[
\mathbf{x}^{(20)}_j
=
\mathbf{x}^{(10)}_{2j}
+
\mathbf{x}^{(10)}_{2j+1},
\]

where the addition is performed separately for TCP, UDP, SSDP, and ARP.  A
30-minute row is the sum of three adjacent ten-minute rows, and a 60-minute row
is the sum of six.

This gives the same category counts that would be obtained by reading the PCAP
again, provided that all of the following remain unchanged:

\begin{itemize}
  \item the device filter;
  \item the four category definitions;
  \item the UTC epoch boundary rule;
  \item the source ten-minute rows are complete and consecutive.
\end{itemize}

No packet download or packet classification is required.  Only the small
processed JSON files are read.

\subsection{Durations that can and cannot be recovered}

\begin{table}[htbp]
\centering
\small
\begin{tabular}{p{0.19\textwidth}p{0.21\textwidth}p{0.48\textwidth}}
\toprule
Requested duration & Recoverable from saved rows? & Reason \\
\midrule
5 minutes
& No
& A ten-minute count does not reveal which packets occurred in its first or
  second half. \\
\addlinespace
10 minutes
& Already available
& This is the stored resolution. \\
\addlinespace
15 minutes
& Not exactly
& A 15-minute boundary cuts some stored ten-minute rows in half. \\
\addlinespace
20 minutes
& Yes
& Sum two correctly aligned ten-minute rows. \\
\addlinespace
30 minutes
& Yes
& Sum three correctly aligned ten-minute rows. \\
\addlinespace
40 or 60 minutes
& Yes
& Sum four or six correctly aligned ten-minute rows. \\
\addlinespace
Any multiple of 10 minutes
& Yes
& The original count intervals can be combined without splitting one. \\
\bottomrule
\end{tabular}
\end{table}

Smaller windows, non-multiple durations, packet directions, byte totals, TCP
flags, and other information not stored in the four counts cannot be recovered
from these JSON files.  The PCAP files would be required for those changes.

\subsection{Correct boundary handling}

Rows must be grouped according to their timestamps and the requested epoch
boundary, not merely taken in pairs or triples starting with the first row in a
file.  This matters when a capture starts partway through a larger interval.

For a 20-minute experiment, for example, the valid groups are the epoch-aligned
intervals ending at minutes 00, 20, and 40 within each hour.  An incomplete
group at the beginning or end of a capture should be dropped or explicitly
marked incomplete; it must not be presented as a complete 20-minute row.

The conversion should also remain inside each capture and research partition.
It must not join the final row of a fitting capture to the first row of a
calibration capture.

\subsection{Silent and missing source rows}

The safest state rules for a longer interval are:

\begin{itemize}
  \item if any required ten-minute row is missing, mark the longer interval as
        missing;
  \item if every required row is observed-silent, mark the longer interval as
        observed-silent;
  \item otherwise, sum the four counts and mark the interval as observed.
\end{itemize}

Treating a missing row as four zeros would confuse unknown capture coverage
with a genuinely silent device and could change both fitting and anomaly
decisions.

\subsection{Reconstructing the attack labels}

The existing attack-label JSON files can also be combined without the PCAPs.
For each longer interval:

\begin{itemize}
  \item set \texttt{is\_attack} to true if any constituent row is attack
        positive;
  \item take the union of the attack types and affected features;
  \item add the non-overlapping attack-overlap seconds from the constituent
        rows.
\end{itemize}

This preserves the current rule that any attack overlap makes the complete row
positive.  It also exposes an evaluation problem: as windows become longer,
more normal minutes are included in positive rows.  Window-level precision,
recall, and F1 therefore cannot be interpreted without also reporting attack
episodes, time to first alert, and false alerts per day.

\section{What must be rerun after reconstructing longer windows}

The expensive packet-reading stage does not need to be repeated, but the
statistical stages do.  A 20-minute experiment requires:

\begin{enumerate}
  \item derive 20-minute count and label files from the saved ten-minute files;
  \item fit new alpha parameters using the 20-minute fitting rows;
  \item calculate a new threshold using the 20-minute calibration rows;
  \item score the 20-minute development rows;
  \item evaluate the new decisions using the reconstructed 20-minute labels.
\end{enumerate}

The ten-minute alpha parameters and threshold must not be reused.  Raw anomaly
scores from different durations are also not directly comparable because the
number of packets and the probability of each row change with duration.

Changing only \texttt{window\_minutes} in the current configuration is not
enough.  That setting controls aggregation while reading PCAPs; it does not
automatically combine already processed JSON rows.  The repository currently
has no processed-row conversion command, so this path would require a small,
validated conversion step.

Derived files and artifacts should be written to new duration-specific paths.
The ten-minute baseline should not be overwritten.  The record for a derived
dataset should include the checksums of its ten-minute source files, the new
duration, the boundary rule, and the conversion code version.

\section{Current data limits for longer-window experiments}

The amount of calibration data places a practical upper limit on window
length.  The current files contain:

\begin{itemize}
  \item 576 ten-minute Chromecast calibration rows;
  \item 81 ten-minute Samsung-camera calibration rows.
\end{itemize}

With non-overlapping, complete windows, the Samsung calibration set would
contain 40 complete 20-minute rows and 27 complete 30-minute rows.  The current
configuration requires at least 30 calibration rows.  Therefore:

\begin{itemize}
  \item a 20-minute comparison is technically possible for both labelled
        devices, although 40 Samsung rows provide a weak basis for estimating a
        first-percentile threshold;
  \item a non-overlapping 30-minute Samsung experiment would fail the current
        minimum-sample requirement;
  \item 60-minute windows would leave still fewer calibration observations and
        should not be treated as a comparable experiment with the present
        partition.
\end{itemize}

Lowering the minimum merely to make a longer-window experiment run would not
create the missing calibration evidence.  A larger benign calibration period
would be the scientifically sound remedy.

\section{Other no-PCAP adjustments}

\subsection{Overlapping longer windows}

A 20-minute total can be calculated every ten minutes by adding the current row
and the previous row.  This provides 20 minutes of context while retaining a
new decision every ten minutes.  The same idea applies to 30-minute totals.

The cost is that neighboring observations share most of their packets.  Their
scores are therefore strongly related, and one attack may produce several
nearly duplicate alerts.  Fitting, calibration, labeling, and episode-level
evaluation must all use the same overlapping policy.  Overlapping rows increase
the row count but do not create additional independent calibration evidence.

\subsection{Rules over the existing ten-minute scores}

The model can remain unchanged while its sequence of scores is interpreted over
time.  Examples include:

\begin{itemize}
  \item open an incident only after two consecutive anomalous rows;
  \item alert when at least two of the last three rows are anomalous;
  \item calculate a running weighted average that gives the greatest weight to
        the newest score.
\end{itemize}

These rules need no count conversion and no PCAPs.  They can reduce isolated
false alerts or reveal a sustained moderate change.  They also delay some
alerts and can miss a single short attack, so the rule must be calibrated on
earlier data and evaluated by attack episode rather than selected after looking
at the final test.

\subsection{Several durations together}

The detector could retain the ten-minute score for short changes and add a
20-minute or 30-minute score for sustained changes.  This avoids choosing one
duration for every attack type, but it creates a new decision rule: for example,
whether either score or both scores must cross their thresholds.  That rule
introduces another opportunity for false alerts and requires its own
calibration.  It should be considered only after the single-duration comparison
is complete.

\section{Recommendation for this project}

Keep the current ten-minute system as the reference result.  The smallest fair
window-length experiment that does not require PCAP files is:

\begin{enumerate}
  \item construct non-overlapping 20-minute rows and labels from pairs of
        epoch-aligned ten-minute rows;
  \item preserve missing intervals and discard incomplete boundary groups;
  \item fit, calibrate, and score a completely separate 20-minute model;
  \item keep the chronological partitions unchanged;
  \item express adaptation intervals in hours before converting them back to a
        number of windows;
  \item compare the 10-minute and 20-minute systems using window metrics,
        attack-episode detection, time to first alert, and false alerts per day.
\end{enumerate}

Twenty minutes is the only larger non-overlapping candidate that both labelled
devices can currently calibrate without violating the configured minimum of 30
rows.  Even then, the Samsung first-percentile threshold will be based on only
40 rows and must be reported as a limitation.

If the aim is to reduce isolated false alarms rather than change the underlying
count model, a two-consecutive-window incident rule is the smaller experiment.
It uses the existing ten-minute scores and requires no count reconstruction.
It should not replace the 20-minute comparison because the two approaches
answer different questions.

\section{Direct answer about alpha aggregation}

The production code does not average alpha parameters from individual rows.
It estimates one set of four alpha parameters from all normal training rows.
It keeps the rows separate while calculating their probabilities and adds the
resulting log-probabilities because that is the logarithm of the probability of
the complete collection.

This calculation retains the usual category composition and the observed
amount of window-to-window variation.  It does not retain time of day, row
order, duration of an activity, or day-specific regimes.  Those temporal
patterns can disappear, but their disappearance is caused by the
time-independent model, not by the mathematical addition of the row
log-probabilities.

\section{Relevant implementation locations}

The behavior described above can be checked directly in the following files:

\begin{itemize}
  \item \path{src/lm_idnet/models/modeling_stage.py}: loads every non-missing
        training window into one count matrix and calls the estimator once;
  \item \path{src/lm_idnet/algorithms/dirichlet.py}: creates one initial set of
        alpha parameters and repeatedly updates that same set using all rows;
  \item \path{src/lm_idnet/algorithms/log_likelihood.py}: calculates a term for
        each row and then sums those row terms;
  \item \path{src/lm_idnet/processing/statistics.py}: optionally estimates
        capture-level or daily values for reporting, separately from production
        model training;
  \item \path{src/lm_idnet/processing/packet_transformer.py}: creates the
        epoch-aligned count rows from packet timestamps;
  \item \path{src/lm_idnet/processing/window_policy.py}: defines the common
        boundary convention;
  \item \path{src/lm_idnet/processing/schemas.py}: preserves timestamps, counts,
        silent intervals, and missing intervals in the processed JSON files.
\end{itemize}

\end{document}
